\documentclass[a4paper,twoside]{article}

\usepackage{morewrites}

\usepackage[svgnames,dvipsnames,x11names]{xcolor}
\usepackage{graphicx}
\usepackage[english]{babel}
\usepackage[colorlinks=true, linkcolor=NavyBlue, urlcolor=NavyBlue, citecolor=NavyBlue]{hyperref}
\usepackage[a4paper]{geometry}
\usepackage{ts-fonts}
\usepackage{amsmath}
\usepackage{float}
\usepackage{seqsplit}
\usepackage{ts-code}

\usepackage{lastpage}
\usepackage{refcount}
\makeatletter
\def\ps@plain{
  \let\@oddhead\@empty
  \let\@evenhead\@empty
  \def\@oddfoot{
    \hfil
    \ifnum\getpagerefnumber{LastPage}>1\relax
      \thepage
    \fi
    \hfil}
  \let\@evenfoot\@oddfoot
}
\makeatother

\title{Progress Bars}
\author{}\date{}

\begin{document}

\maketitle

\thispagestyle{plain}
The \texttt{texsmith.progressbar} extension mirrors the \href{https://facelessuser.github.io/pymdown-extensions/extensions/progressbar/}{\texttt{pymdownx.progressbar}} syntax and
converts it into \LaTeX{} commands powered by the
\href{https://ctan.org/pkg/progressbar}{\texttt{progressbar}} package. It is enabled by default, so you can
drop the \texttt{[=75\% "Done"]} shorthand in your Markdown and keep both the HTML preview and the PDF build in sync.

\section{Syntax}\label{syntax}

\begin{code}{markdown}{}{}
\begin{Verbatim}[commandchars=\\\{\},breaklines, breakanywhere, commandchars=\\\{\}]
[=25\PYZpc{} \PYZdq{}Research\PYZdq{}]
[=50\PYZpc{} \PYZdq{}Implementation\PYZdq{}]
[=75\PYZpc{} \PYZdq{}Review\PYZdq{}]
[=100\PYZpc{} \PYZdq{}Launch\PYZdq{}]\PYZob{}: .thin\PYZcb{}
\end{Verbatim}
\end{code}
\begin{itemize}
\item{} Values must be percentages (\texttt{0 – 100}). TeXSmith clamps the values if needed.
\item{} The quoted label is optional; when omitted the percentage is used.
\item{} Trailing attribute lists (\texttt{\{: .class \#id \}}) attach CSS classes and HTML attributes.
  Use the \texttt{.thin} class to halve the bar height (e.g. for tables or dense summaries).

\end{itemize}
\section{LaTeX output}\label{latex-output}

During rendering, TeXSmith emits a \texttt{\textbackslash{}progressbar} call with the following defaults:

\begin{code}{latex}{}{}
\begin{Verbatim}[commandchars=\\\{\},breaklines, breakanywhere, commandchars=\\\{\}]
\PY{n+nb}{\PYZob{}}\PY{k}{\PYZbs{}progressbar}[
  width=9cm,
  heighta=12pt,
  roundnessr=0.1,
  borderwidth=1pt,
  linecolor=black,
  filledcolor=black!60,
  emptycolor=black!10
]\PY{n+nb}{\PYZob{}}0.73\PY{n+nb}{\PYZcb{}} Launch\PY{n+nb}{\PYZcb{}}
\end{Verbatim}
\end{code}
The \texttt{.thin} class switches \texttt{heighta} to \texttt{6pt}. If you need more control, wrap the Markdown block
in a raw \LaTeX{} fence and tweak the options directly.

\section{Example project}\label{example-project}

Use the bundled \texttt{examples/progressbar} project for smoke tests or screenshots:

\begin{code}{bash}{}{}
\begin{Verbatim}[commandchars=\\\{\},breaklines, breakanywhere, commandchars=\\\{\}]
\PY{n+nb}{cd}\PY{+w}{ }examples/progressbar
texsmith\PY{+w}{ }progressbar.md\PY{+w}{ }\PYZhy{}\PYZhy{}template\PY{+w}{ }article\PY{+w}{ }\PYZhy{}\PYZhy{}output\PYZhy{}dir\PY{+w}{ }build\PY{+w}{ }\PYZhy{}\PYZhy{}build
\end{Verbatim}
\end{code}
\begin{figure}[H]
    \centering

    \ifdefined\adjustbox

        \adjustbox{max width=\textwidth}{
            \includegraphics[width=\linewidth]{snippet-<HASH>.pdf}
        }

    \else

        \includegraphics[width=\linewidth]{snippet-<HASH>.pdf}

    \fi

\end{figure}

\end{document}
